\setvseed{\VAR{seed}}
\skillheader[Fundamentals Checkpoint]{FCP}

\#{ Every student's paper sits in one document, so the problem counter has to
   start over here or the second student's sheet is numbered 21 to 40. }
\setcounter{pi}{0}

\textit{Calculate each of the following. You do not need to simplify any fractions.}

\vspace{6pt}

\begin{tabular}{lp{4.6cm}lp{4.6cm}lp{4.6cm}}
    \numprob
    $\VAR{p1_prob}$ \newline\newline\ans{$\VAR{p1_ans}$} &

    \numprob
    $\VAR{p2_prob}$ \newline\newline\ans{$\VAR{p2_ans}$} &

    \numprob
    $\VAR{p3_prob}$ \newline\newline\ans{$\VAR{p3_ans}$} \\[1.7in]

    \numprob
    $\VAR{p4_prob}$ \newline\newline\ans{$\VAR{p4_ans}$} &

    \numprob
    $\VAR{p5_prob}$ \newline\newline\ans{$\VAR{p5_ans}$} &

    \numprob
    $\VAR{p6_prob}$ \newline\newline\ans{$\VAR{p6_ans}$} \\[1.7in]

    \numprob
    \begin{tabular}[t]{llr}
\BLOCK{for part in p7_parts}
        (\VAR{part.letter}) & $\VAR{part.prob}$ & \ans{$\VAR{part.ans}$} \BLOCK{if not loop.last}\\[12pt]\BLOCK{endif}
\BLOCK{endfor}
    \end{tabular} &

    \numprob
    \begin{tabular}[t]{llr}
\BLOCK{for part in p8_parts}
        (\VAR{part.letter}) & $\VAR{part.prob}$ & \ans{$\VAR{part.ans}$} \BLOCK{if not loop.last}\\[12pt]\BLOCK{endif}
\BLOCK{endfor}
    \end{tabular} &

    \numprob
    \begin{tabular}[t]{llr}
\BLOCK{for part in p9_parts}
        (\VAR{part.letter}) & $\VAR{part.prob}$ & \ans{$\VAR{part.ans}$} \BLOCK{if not loop.last}\\[12pt]\BLOCK{endif}
\BLOCK{endfor}
    \end{tabular} \\[1.7in]

    \numprob
    \begin{tabular}[t]{llr}
\BLOCK{for part in p10_parts}
        (\VAR{part.letter}) & $\VAR{part.prob}$ & \ans{$\VAR{part.ans}$} \BLOCK{if not loop.last}\\[12pt]\BLOCK{endif}
\BLOCK{endfor}
    \end{tabular} &

    \numprob
    $\VAR{p11_prob}$ \newline\newline\newline\ans{ $\VAR{p11_ans}$ } &

    \numprob
    $\VAR{p12_prob}$ \newline\newline\newline\ans{ $\VAR{p12_ans}$ }
\end{tabular}

\newpage
\#{ The version and the grading box belong on the first sheet only. Set for
   this page alone, never with \fancyfoot, because every student's paper is in
   one document and a running change would follow through to all of them. }
\thispagestyle{skillcontinued}

\#{ A line of air under the header. Page one has the skill box between the
   header and the first problem; page two has nothing, so without this the top
   row sits right under the student's name. \vspace* rather than \vspace,
   because ordinary vertical space is discarded at the top of a page. }
\vspace*{\baselineskip}

\begin{tabular}{lp{7.9cm}lp{7.9cm}}
    \numprob
    $\VAR{p13_prob}$ \newline\newline\newline\ans{ $\VAR{p13_ans}$ } &

    \numprob
    $\VAR{p14_prob}$ \newline\newline\newline\ans{ $\VAR{p14_ans}$ } \\[2in]

    \numprob
    $\VAR{p15_prob}$ \newline\newline\ans{$\VAR{p15_ans}$} &

    \numprob
    $\VAR{p16_prob}$ \newline\newline\ans{$\VAR{p16_ans}$} \\[2in]
\end{tabular}

\textit{For each of the following, order the given integers from least to greatest.}

\begin{tabular}{lp{7.9cm}lp{7.9cm}}
    \numprob
    $\VAR{p17_prob}$ \newline\newline \ans{$\VAR{p17_ans}$} &

    \numprob
    \#{ Half the space the other blocks get. Ordering six integers is a single
       line of writing, not worked steps, so 2in of it was mostly empty paper.
       The slack also pays for the line of air under the header above. }
    $\VAR{p18_prob}$ \newline\newline \ans{$\VAR{p18_ans}$} \\[0.9in]
\end{tabular}

\textit{Round each of the following to the desired place value.}

\begin{tabular}{lp{7.9cm}lp{7.9cm}}
    \numprob
    \begin{tabular}[t]{lp{7cm}}
\BLOCK{for part in p19_parts}
        (\VAR{part.letter}) & $\VAR{part.prob}$ (\VAR{part.place}) \newline\newline\ans{$\VAR{part.ans}$} \BLOCK{if not loop.last}\\[48pt]\BLOCK{endif}
\BLOCK{endfor}
    \end{tabular} &

    \numprob
    \begin{tabular}[t]{lp{7cm}}
\BLOCK{for part in p20_parts}
        (\VAR{part.letter}) & $\VAR{part.prob}$ (\VAR{part.place}) \newline\newline\ans{$\VAR{part.ans}$} \BLOCK{if not loop.last}\\[48pt]\BLOCK{endif}
\BLOCK{endfor}
    \end{tabular}
\end{tabular}
